% 02-modele.tex — Les objets simulés (barre, corde, seau, eau) et les choix qui
% demandent une explication.

\section{Les objets simulés}

La scène contient quatre objets (\cref{fig-schema}) ; hormis son poids, toutes
les forces sur le seau passent par la corde.
\begin{itemize}
  \item \emph{La barre} (\SI{0,9}{m}, \SI{0,3}{kg}) tourne sur un pivot sans
        frottement. Pour la rotation, une barre homogène de moment d'inertie
        \(m_b L_b^2/3\) équivaut à un point de masse \(m_b/3\) à son bout ; pour
        que le moment de son poids, \(m_b g (L_b/2) \sin(\alpha)\), soit juste,
        ce point subit une force \(m_b g/2\). Tant que le moteur la mène, sa
        position est imposée.
  \item \emph{La corde} (\SI{0,15}{kg/m}) est une chaîne de dix maillons qui
        peuvent raccourcir mais pas s'allonger : elle plie, se détend, ne pousse
        jamais. Son dernier nœud est l'anse du seau.
  \item \emph{Le seau} est un cylindre rigide (Ø \SI{26}{cm}, \SI{24}{cm} de
        profondeur, \SI{0,45}{kg}) avec son moment d'inertie (paroi et fond),
        pendu à une anse de \SI{0,05}{kg} qui pivote aux oreilles ; un petit
        frottement amortit ce pivot et le bout de la barre (\SI{8}{s^{-1}}).
  \item \emph{L'eau}, \(m - \SI{0,5}{kg}\) à \SI{1}{kg} par litre, est faite de
        dix billes par défaut, qui portent chacune un dixième de cette masse,
        roulent sur les parois, se poussent et poussent le seau. L'essai est
        raté dès qu'une bille franchit le bord.
\end{itemize}
Comme l'eau a une masse et une place, le centre de masse du seau rempli dépend
de la quantité d'eau. \(R\) est donc la distance du pivot au centre de masse
calculé sur les billes au repos : pour \(L = \SI{0,4}{m}\), \(R\) vaut \SI{1,585}{m} à \SI{0,6}{kg},
\SI{1,650}{m} à \SI{2}{kg} et \SI{1,622}{m} à \SI{8}{kg}
(\cref{tab-repos}). C'est ce \(R\) qui sert dans \(\sqrt{5 g R}\).

\begin{figure}[ht]
  \centering
  \begin{tikzpicture}[scale=1.05, every node/.style={font=\footnotesize}]
    \coordinate (pivot) at (0,0);
    \coordinate (tip) at (2.4,-1.2);
    \coordinate (anse) at (4.6,-2.45);
    \draw[metal!40, line width=3pt] (-0.5,-3.1) -- (pivot) -- (0.5,-3.1);
    \draw[metal, line width=2.6pt] (pivot) -- (tip);
    \fill (pivot) circle (2.4pt);
    \node[above left] at (pivot) {pivot};
    \node[above=2pt, sloped, metal] at ($(pivot)!0.35!(tip)$) {barre \(L_b\), \(m_b\)};
    \foreach \k in {1,...,9} { \coordinate (r\k) at ($(tip)!\k/10!(anse)+(0,-0.028*\k+0.0028*\k*\k)$); }
    \draw[corde, line width=1.1pt] (tip) \foreach \k in {1,...,9} { -- (r\k) } -- (anse);
    \foreach \k in {1,...,9} { \fill[corde] (r\k) circle (1.3pt); }
    \fill (tip) circle (1.8pt);
    \node[corde, above right] at ($(r4)+(0,0.08)$) {10 maillons \(\ell = L/10\), masse \(\mu \ell\) par nœud};
    \begin{scope}[shift={(anse)}, rotate=-30]
      \coordinate (ears) at (0,-0.42);
      \draw[black] (-0.34,-0.42) .. controls (-0.3,0.05) and (0.3,0.05) .. (0.34,-0.42);
      \draw[metal, very thick, fill=metal!10] (-0.34,-0.42) -- (-0.3,-1.05) -- (0.3,-1.05) -- (0.34,-0.42);
      \foreach \x in {-0.24,-0.12,0,0.12,0.24} { \fill[eau] (\x,-0.98) circle (2pt); }
      \foreach \x in {-0.18,-0.06,0.06,0.18} { \fill[eau] (\x,-0.88) circle (2pt); }
      \fill[metal] (-0.34,-0.42) circle (1.3pt) (0.34,-0.42) circle (1.3pt);
      \node[left=4pt] at (-0.34,-0.42) {oreilles};
      \coordinate (cm) at (0,-0.82);
      \fill (cm) circle (1.6pt);
    \end{scope}
    \fill (anse) circle (1.8pt);
    \node[above right] at (anse) {anse (nœud 10)};
    \draw[-{Stealth}, poids, thick] (cm) -- ++(0,-0.7) node[right] {\(\vect{P} = m \vect{g}\)};
    \draw[-{Stealth}, tensionc, thick] (anse) -- ($(anse)!0.7cm!(r9)$) node[above left=-1pt] {\(\vect{T}\)};
    \draw[gray, dashed] (pivot) -- (cm);
    \node[gray, below left] at ($(pivot)!0.62!(cm)$) {\(R\)};
    \node[eau, right] at ($(cm)+(0.45,-0.05)$) {billes d'eau};
  \end{tikzpicture}
  \caption{Modèle simulé : la barre, les dix maillons de la corde (l'anse est le
    dixième nœud), le seau rigide articulé aux oreilles et les billes d'eau.
    \(R\) va du pivot au centre de masse du seau rempli (point noir). Seuls le
    poids \(\vect{P}\) et la tension \(\vect{T}\) de la corde agissent sur le
    seau rempli.}
  \label{fig-schema}
\end{figure}
